\documentclass[11pt,letterpaper]{article}

\usepackage[utf8]{inputenc}
\usepackage[T1]{fontenc}
\usepackage{lmodern}
\usepackage[spanish,es-noshorthands,es-tabla,es-nodecimaldot]{babel}
\usepackage[margin=2.5cm]{geometry}
\usepackage{amsmath,amssymb}
\usepackage{graphicx}
\usepackage{subcaption}
\usepackage{booktabs}
\usepackage{float}
\usepackage{xcolor}
\usepackage{listings}
\usepackage[hidelinks]{hyperref}

% ---------- Formato del código ----------
\definecolor{fondo}{RGB}{247,247,247}
\definecolor{comentario}{RGB}{95,125,95}
\definecolor{clave}{RGB}{0,70,150}
\definecolor{cadena}{RGB}{160,60,40}

\lstset{
    language=Python,
    basicstyle=\ttfamily\small,
    keywordstyle=\color{clave}\bfseries,
    commentstyle=\color{comentario}\itshape,
    stringstyle=\color{cadena},
    backgroundcolor=\color{fondo},
    frame=single,
    rulecolor=\color{black!25},
    showstringspaces=false,
    breaklines=true,
    columns=fullflexible,
    keepspaces=true,
    captionpos=b,
    literate={á}{{\'a}}1 {é}{{\'e}}1 {í}{{\'i}}1 {ó}{{\'o}}1 {ú}{{\'u}}1 {ñ}{{\~n}}1
}
\renewcommand{\lstlistingname}{Código}
\setlength{\emergencystretch}{3em}  % evita líneas desbordadas por nombres en \texttt

% ---------- Figura con un grafo generado y su árbol de Dijkstra ----------
% #1 = nombre base del archivo (sin extensión), #2 = título, #3 = etiqueta
\newcommand{\grafoyarbol}[3]{%
\begin{figure}[H]
    \centering
    \begin{subfigure}{\textwidth}
        \centering
        \includegraphics[width=\linewidth,height=0.38\textheight,keepaspectratio]{imagenes/generados/#1.png}
        \caption{Grafo generado}
    \end{subfigure}

    \medskip
    \begin{subfigure}{\textwidth}
        \centering
        \includegraphics[width=\linewidth,height=0.42\textheight,keepaspectratio]{imagenes/calculados/#1_Dijkstra.png}
        \caption{Árbol de caminos más cortos}
    \end{subfigure}
    \caption{#2. El nodo fuente aparece en rojo.}
    \label{#3}
\end{figure}
\clearpage
}

\title{Proyecto 3\\Algoritmo de Dijkstra}
\author{Josue gutierrez Andrade\\[4pt]\small Diseño y Análisis de Algoritmos}
\date{}

\begin{document}

\maketitle

\begin{abstract}
Sobre la biblioteca de grafos de los proyectos anteriores se implementa el método
\texttt{Dijkstra} de la clase \texttt{Grafo}. Dado un nodo fuente $s$, el método
calcula el árbol de caminos más cortos y nombra cada nodo con su distancia a $s$;
por ejemplo, \texttt{"2 (22.45)"}. El método se aplicó a 10 grafos con pesos
aleatorios, generados con los cinco modelos aleatorios de la biblioteca con 30 y
500 nodos. Todos los grafos y árboles se exportaron a formato GraphViz y se
visualizaron.
\end{abstract}

\tableofcontents
\clearpage

% =====================================================================
\section{Introducción}

En un grafo con pesos no negativos en las aristas, la longitud de un camino es la
suma de los pesos de sus aristas. El problema del camino más corto consiste en
encontrar, para cada nodo $v$, un camino de longitud mínima desde un nodo fuente
$s$. Esa longitud mínima es la distancia $d(v)$. El algoritmo de Dijkstra resuelve
el problema para todos los nodos a la vez. Si para cada nodo $v \neq s$ se guarda
el nodo desde el que se llegó a él por el camino más corto, esas aristas forman un
árbol con raíz en $s$: el \emph{árbol de caminos más cortos}. En él, el único
camino de $s$ a $v$ es un camino más corto en el grafo original.

% =====================================================================
\section{Pesos en la biblioteca de grafos}

Los grafos de los proyectos anteriores no tenían pesos, así que se hicieron estos
cambios en la clase \texttt{Grafo}:

\begin{itemize}
    \item \texttt{agregar\_arista(u, v, peso)} guarda el peso en un diccionario
          \texttt{\_pesos}, indexado por el par $(u, v)$; en un grafo no dirigido,
          también por $(v, u)$. Con él, \texttt{peso(u, v)} consulta el peso de una
          arista en tiempo $O(1)$.
    \item \texttt{asignar\_pesos\_aleatorios(minimo, maximo)} asigna a cada arista
          un peso uniforme en $[\text{minimo}, \text{maximo}]$, redondeado a dos
          decimales.
    \item \texttt{guardar\_graphviz} escribe el peso en el atributo
          \texttt{weight}; por ejemplo, \texttt{"0" -{}- "26" [weight=11.07];}.
\end{itemize}

Los métodos \texttt{BFS}, \texttt{DFS\_R} y \texttt{DFS\_I} del Proyecto~2 se
conservan sin cambios.

% =====================================================================
\section{Algoritmo de Dijkstra}

\subsection{Descripción}

Se sigue el algoritmo visto en clase. Sea $S$ el conjunto de nodos explorados,
para los cuales ya se conoce la distancia definitiva. Al inicio $S = \emptyset$,
$d(s) = 0$ y la distancia tentativa de los demás nodos es $\infty$. En cada paso:

\begin{enumerate}
    \item Se elige el nodo $u \notin S$ con la menor distancia tentativa
          (\emph{delmin}) y se agrega a $S$.
    \item Para cada arista $(u, v)$ con $v \notin S$, si
          $d(u) + \ell(u, v) < d(v)$, se actualiza $d(v) \leftarrow d(u) + \ell(u, v)$
          y el padre de $v$ pasa a ser $u$ (\emph{decrease key}).
\end{enumerate}

El proceso termina cuando ya no quedan nodos alcanzables fuera de $S$. Las aristas
$(\text{padre}(v), v)$ forman el árbol de caminos más cortos.

\subsection{Implementación}

La cola de prioridad es un montículo binario (módulo \texttt{heapq}). Como
\texttt{heapq} no permite disminuir la prioridad de un elemento, el \emph{decrease
key} se implementa insertando una nueva entrada con la distancia mejorada. Las
entradas viejas de un nodo quedan en el montículo y se descartan al extraerlas si
el nodo ya está en $S$. Cada entrada lleva además un contador de desempate, para
que el montículo nunca tenga que comparar dos identificadores de nodo.

\begin{lstlisting}[caption={Método \texttt{Dijkstra}.}]
def Dijkstra(self, s):
    self._validar_fuente(s)
    for arista in self.aristas:
        if arista.peso is not None and arista.peso < 0:
            raise ValueError("Dijkstra no admite aristas con peso negativo")

    distancia = {s: 0}
    padre = {s: None}
    explorados = []         # nodos en el orden en que entran a S
    en_s = set()
    desempate = itertools.count()   # evita comparar identificadores de nodos
    cola = [(0, next(desempate), s)]
    while cola:
        d_u, _, u = heapq.heappop(cola)          # delmin
        if u in en_s:
            continue                             # entrada obsoleta
        en_s.add(u)
        explorados.append(u)
        for v in self._vecinos[u]:
            if v in en_s:
                continue
            peso = self._pesos[(u, v)]
            nueva = d_u + (1 if peso is None else peso)
            if v not in distancia or nueva < distancia[v]:
                distancia[v] = nueva             # decrease key
                padre[v] = u
                heapq.heappush(cola, (nueva, next(desempate), v))

    arbol = Grafo(self.dirigido)
    etiqueta = {v: f"{v} ({distancia[v]:.2f})" for v in explorados}
    for v in explorados:
        arbol.agregar_nodo(etiqueta[v])
        if padre[v] is not None:
            arbol.agregar_arista(etiqueta[padre[v]], etiqueta[v], self._pesos[(padre[v], v)])
    return arbol
\end{lstlisting}

El árbol resultante es un nuevo \texttt{Grafo}. Cada nodo se llama
\texttt{"nodo (distancia)"}, con la distancia redondeada a dos decimales, y cada
arista conserva su peso. El árbol contiene solo los nodos alcanzables desde $s$.
Si hay pesos negativos el método lanza \texttt{ValueError}, porque Dijkstra no es
correcto en ese caso, y si $s$ no existe lanza \texttt{KeyError}. Las aristas sin
peso cuentan como peso 1, en cuyo caso las distancias coinciden con las de BFS.

\subsection{Complejidad y correctitud}

Cada arista puede provocar a lo más una inserción en el montículo, así que este
recibe $O(m)$ entradas, y cada inserción o extracción cuesta $O(\log m) = O(\log n)$.
El costo total es $O((n + m) \log n)$.

La correctitud se basa en el teorema visto en clase: cuando un nodo $v$ entra a
$S$, su distancia tentativa ya es la distancia mínima. Si existiera un camino $P$
más corto, tendría una primera arista $(x, y)$ que sale de $S$, y como los pesos
son no negativos, $L(P) \geq d(x) + \ell(x, y) \geq d(y) \geq d(v)$, porque $v$ se
eligió con la menor distancia tentativa. Eso contradice que $P$ sea más corto.

\subsection{Ejemplo}

Como prueba se usó el grafo dirigido del ejemplo de clase, con origen $s$. La
tabla~\ref{tab:ejemplo} muestra el orden en que los nodos entran a $S$ y la
etiqueta que reciben en el árbol; coincide con el resultado de las diapositivas.

\begin{table}[H]
\centering
\begin{tabular}{clr}
\toprule
Orden & Nodo en el árbol & Padre \\
\midrule
1 & \texttt{s (0.00)} & --- \\
2 & \texttt{2 (9.00)} & $s$ \\
3 & \texttt{6 (14.00)} & $s$ \\
4 & \texttt{7 (15.00)} & $s$ \\
5 & \texttt{3 (32.00)} & 6 \\
6 & \texttt{5 (34.00)} & 3 \\
7 & \texttt{4 (45.00)} & 5 \\
8 & \texttt{t (50.00)} & 5 \\
\bottomrule
\end{tabular}
\caption{Resultado de \texttt{Dijkstra("s")} en el ejemplo de clase.}
\label{tab:ejemplo}
\end{table}

\subsection{Verificación}

Además del ejemplo anterior, se comprobó con un programa de prueba que, en grafos
de los cinco modelos aleatorios con 30 y 500 nodos y pesos aleatorios, se cumple
lo siguiente:

\begin{itemize}
    \item La distancia de cada nodo coincide con la calculada de forma independiente
          por el algoritmo de Bellman--Ford.
    \item Cada etiqueta tiene el formato \texttt{"nodo (distancia)"}, con dos
          decimales.
    \item El árbol tiene un nodo por cada nodo alcanzable, un solo padre por nodo,
          $n - 1$ aristas, y cada arista conserva el peso del grafo original.
    \item Para cada arista $(u, v)$ del árbol, $d(v) = d(u) + \ell(u, v)$; es decir,
          todas las aristas del árbol están en caminos más cortos.
    \item El método también funciona en grafos dirigidos, sin pesos o no conexos, y
          rechaza pesos negativos y nodos fuente inexistentes.
\end{itemize}

Todo el código, incluidos los scripts de generación y de imágenes, se ejecutó con
Python~3.6.8, la versión solicitada.

% =====================================================================
\section{Experimentos}

El script \texttt{generar\_gv.py} genera dos grafos por cada uno de los cinco
generadores aleatorios: uno con ``pocos'' nodos (30) y otro con ``muchos'' (500).
La malla no se usa porque no es aleatoria. Cada grafo se vuelve a generar hasta
que resulta conexo, de modo que el árbol abarca todos sus nodos. Después se le
asignan pesos aleatorios uniformes entre 1 y 100 y se calcula su árbol de caminos
más cortos desde el nodo $0$.

\begin{table}[H]
\centering
\begin{tabular}{ll}
\toprule
Modelo & Parámetros \\
\midrule
Erd\H{o}s--Rényi & $m = \operatorname{round}\!\left(\dfrac{n(\ln n + 2)}{2}\right)$ \\[10pt]
Gilbert & $p = \dfrac{\ln n + 3}{n}$ \\[10pt]
Geográfico simple & $r = \sqrt{\dfrac{\ln n + 4}{\pi n}}$ \\[10pt]
Barabási--Albert & $d = 4$ \\[4pt]
Dorogovtsev--Mendes & sin parámetros adicionales \\
\bottomrule
\end{tabular}
\caption{Parámetros de los generadores.}
\end{table}

Los grafos se guardan en \texttt{gv/generados/} y los árboles en
\texttt{gv/calculados/}. El script \texttt{generar\_imagenes.py} los dibuja con
Graphviz. Los grafos generados se dibujan con \texttt{sfdp}, y en los de 30 nodos
se escribe el peso sobre cada arista. Los árboles se dibujan con \texttt{twopi},
con la raíz al centro y la etiqueta \texttt{nodo (distancia)} en cada nodo.

% =====================================================================
\section{Resultados}

La tabla~\ref{tab:resultados} resume los árboles obtenidos. La \emph{altura} es el
mayor número de aristas en un camino del árbol desde $s$. Como comparación se
incluye la altura del árbol BFS del mismo grafo, que ignora los pesos. Los grafos y
los pesos son aleatorios, así que los valores corresponden a la ejecución incluida
en el repositorio.

\begin{table}[H]
\centering
\begin{tabular}{lrrrrrr}
\toprule
 & & & \multicolumn{2}{c}{Distancia a $s$} & \multicolumn{2}{c}{Altura} \\
\cmidrule(lr){4-5}\cmidrule(l){6-7}
Grafo & Nodos & Aristas & Máxima & Promedio & Dijkstra & BFS \\
\midrule
Erd\H{o}s--Rényi & 30 & 81 & 122.92 & 63.63 & 6 & 3 \\
Gilbert & 30 & 83 & 117.53 & 55.20 & 4 & 3 \\
Geográfico simple & 30 & 84 & 238.10 & 61.61 & 8 & 8 \\
Barabási--Albert & 30 & 56 & 194.53 & 105.90 & 9 & 6 \\
Dorogovtsev--Mendes & 30 & 57 & 188.03 & 102.24 & 5 & 4 \\
\midrule
Erd\H{o}s--Rényi & 500 & 2054 & 162.04 & 88.63 & 10 & 5 \\
Gilbert & 500 & 2281 & 125.69 & 68.82 & 10 & 4 \\
Geográfico simple & 500 & 2361 & 341.20 & 149.57 & 26 & 17 \\
Barabási--Albert & 500 & 997 & 3390.00 & 1747.04 & 114 & 89 \\
Dorogovtsev--Mendes & 500 & 997 & 300.25 & 116.36 & 9 & 6 \\
\bottomrule
\end{tabular}
\caption{Distancias y alturas de los árboles de caminos más cortos.}
\label{tab:resultados}
\end{table}

De la tabla se desprenden varias observaciones:

\begin{itemize}
    \item \textbf{Dijkstra produce árboles más altos que BFS.} En casi todos los
          grafos, el camino más corto por peso usa más aristas que el camino con
          menos aristas: conviene dar varios pasos por aristas ligeras antes que
          uno solo por una arista pesada. Con 500 nodos, el árbol de
          Erd\H{o}s--Rényi tiene altura 10 con Dijkstra y 5 con BFS.
    \item \textbf{En los grafos densos las distancias apenas crecen.} En
          Erd\H{o}s--Rényi y Gilbert, la distancia máxima pasa de unos 120 con 30
          nodos a entre 125 y 162 con 500 nodos. Con más aristas hay más rutas
          alternativas y es más fácil encontrar una formada por aristas ligeras.
    \item \textbf{Barabási--Albert tiene las distancias más grandes.} La variante
          implementada limita el grado de cada nodo a $d = 4$, lo que produce
          caminos muy largos: con 500 nodos, el nodo más lejano está a 114 aristas
          y a una distancia de 3390.
    \item \textbf{En el modelo geográfico la altura crece con el tamaño.} Como las
          aristas solo unen nodos cercanos, llegar a los nodos lejanos requiere
          muchas aristas, igual que en BFS.
\end{itemize}

% =====================================================================
\section{Visualización}

Las figuras~\ref{fig:er30} a~\ref{fig:dm500} muestran cada grafo generado junto con
su árbol de caminos más cortos. En los árboles de 500 nodos las etiquetas son
pequeñas; se leen con zoom en el PDF o en las imágenes de \texttt{imagenes/calculados/}.
\clearpage

\grafoyarbol{1_Erdos_Renyi_n30}{Erd\H{o}s--Rényi con 30 nodos}{fig:er30}
\grafoyarbol{2_Gilbert_n30}{Gilbert con 30 nodos}{fig:gil30}
\grafoyarbol{3_Geografico_n30}{Geográfico simple con 30 nodos}{fig:geo30}
\grafoyarbol{4_Barabasi_Albert_n30}{Barabási--Albert con 30 nodos}{fig:ba30}
\grafoyarbol{5_Dorogovtsev_Mendes_n30}{Dorogovtsev--Mendes con 30 nodos}{fig:dm30}

\grafoyarbol{1_Erdos_Renyi_n500}{Erd\H{o}s--Rényi con 500 nodos}{fig:er500}
\grafoyarbol{2_Gilbert_n500}{Gilbert con 500 nodos}{fig:gil500}
\grafoyarbol{3_Geografico_n500}{Geográfico simple con 500 nodos}{fig:geo500}
\grafoyarbol{4_Barabasi_Albert_n500}{Barabási--Albert con 500 nodos}{fig:ba500}
\grafoyarbol{5_Dorogovtsev_Mendes_n500}{Dorogovtsev--Mendes con 500 nodos}{fig:dm500}

\end{document}
